\documentclass[10pt,a4paper]{amsart}
\usepackage[margin=19mm]{geometry}
\usepackage{amsmath,amssymb,mathtools}
\usepackage[scr=rsfs,cal=euler]{mathalfa}
\usepackage{xfrac}
\usepackage[unicode,hidelinks,bookmarks=false]{hyperref}
\newcommand{\ie}{i.e.\ }
\newcommand{\cf}{cf.\ }
\newcommand{\resp}{respectively\ }
\newcommand{\Subsection}{\subsection}
\providecommand{\nobreakdash}{\nobreak\hbox{-}}
\setlength{\emergencystretch}{2em}
\multlinegap0pt
\usepackage{bm}
\usepackage{bbm}
\usepackage{dsfont}
\usepackage{xspace}
\usepackage{cancel}
\usepackage{mathtools}
\usepackage{amsthm}
\makeatletter
\@ifundefined{theorem}{\newtheorem{theorem}{Theorem}}{}
\@ifundefined{definition}{\newtheorem{definition}[theorem]{Definition}}{}
\@ifundefined{proposition}{\newtheorem{proposition}[theorem]{Proposition}}{}
\makeatother
\newcommand \C { \mathbb{C} }
\newcommand \R { \mathbb{R} } 
\title[Quantum Measurement Trees, II: Quantum Observables as Ortho-]{Quantum Measurement Trees, II: Quantum Observables as Ortho-Measurable Functions and Density Matrices as Ortho-Probability Measures}
\author{Peter J. Hammond}
\date{Source: arXiv:2509.22617v1 (CC BY 4.0)}
\begin{document}
\maketitle

\noindent\emph{Source and licence.} Quantum Measurement Trees, II: Quantum Observables as Ortho-Measurable Functions and Density Matrices as Ortho-Probability Measures. Version arXiv:2509.22617v1 (CC BY 4.0), distributed under the Creative Commons Attribution 4.0 International licence, as recorded in the arXiv licence metadata for this exact version and shown on the abstract page https://arxiv.org/abs/2509.22617v1. Full rights notice: licenses/arxiv-2509.22617-CC-BY-4.0.txt.\par\smallskip

\noindent\textit{Editorial scope.} Counted unit: the Theorem whose source label is \texttt{thm:threeBijections} in the article (source0.tex lines 801--811, source label \texttt{thm:threeBijections}), delivered together with the article's own complete original proof of it (source0.tex lines 813--859, one \texttt{proof} environment of 47 lines). No mathematics is written, repaired or re-derived: statement and proof are the article's own text line for line. Required context retained uncounted: prop:spectralThm (lines 615--631); eq:spectralDecomp (lines 624--627); def:orthMeasFn (lines 786--797). Every definition, display and label of those lines is retained unchanged. Omitted from this package and not redistributed: 1--614, 628--785, 798--800, 812--812, 860--1567. Nothing inside a retained excerpt is omitted or altered: every retained body is delivered exactly as the article's lines are written, so no omission receipt of any kind accompanies this package. Citations: the retained excerpts carry no citation command at all, so no bibliography entry of the article is transcribed and no reference is redistributed. Reference adaptation: none was needed and none was made; every reference inside the delivered unit resolves inside it, so no unresolved reference or citation is produced. Printed slips kept literal: no printed word, symbol, hypothesis or number of the counted statement or of its proof was altered. Layout and class: the article's own class is \texttt{article} with packages including amsmath, amssymb, amsfonts, amsthm, url, graphicx, enumerate, cases, eucal, xcolor; the wrapper is the lane's pinned \texttt{amsart} editorial wrapper at 10pt on A4, which restates the theorem-like environments the retained text needs and adds the article's own macro definitions that the retained text needs (\texttt{\textbackslash C}, \texttt{\textbackslash R}), with extra wrapper packages (bm, bbm, dsfont, xspace, cancel, mathtools). The delivered numbering is the unit's own. Assets: no retained span contains a figure environment, a graphics-inclusion command, a PDF-page inclusion, a table or a TikZ picture; no image file is copied into this package, the package declares no asset dependency and no image file is redistributed with it. Licence: version arXiv:2509.22617v1 of this article is distributed under the Creative Commons Attribution 4.0 International licence, as recorded in the arXiv licence metadata for this exact version and shown on the abstract page https://arxiv.org/abs/2509.22617v1. A full rights notice is delivered at licenses/arxiv-2509.22617-CC-BY-4.0.txt. Characters: every retained excerpt is pure ASCII, so the delivered engine, XeLaTeX with its default Latin Modern text font, reports no missing character.
\begin{proposition} \label{prop:spectralThm}
	Any $n \times n$ Hermitian matrix $ \mathbf A$:
\begin{enumerate}
	\item induces a numerically identified orthogonal decomposition 
 \( \bigoplus\nolimits _{ \lambda \in s^{ \mathbf A} } L_\lambda \) of $ \C^n $
	into linear subspaces $ L_\lambda = E_\lambda \cup \{ \mathbf 0\} $
	that correspond to the eigenspaces $ E_\lambda $ of $ \mathbf A$,
	one for each $ \lambda \in s^{ \mathbf A} $ in the spectrum of $ \mathbf A$;
	\item has a \emph{spectral decomposition} into the eigenvalue-weighted linear combination 
\begin{equation} \label{eq:spectralDecomp}
	\mathbf A = \sum\nolimits _{ \lambda \in s^{ \mathbf A} } 
	\lambda \, \boldsymbol \Pi _{ L_\lambda }
\end{equation}
	of orthogonal projections $ \boldsymbol \Pi _{ L_\lambda }$ 
	onto the corresponding mutually orthogonal eigenspaces $ L_\lambda $.
\end{enumerate}
\end{proposition}

\begin{definition} \label{def:orthMeasFn}
	A function \( \C^n \owns \mathbf x \mapsto f( \mathbf x) \in \R \cup \{*\} \) 
	is \emph{ortho-measurable} just in case there exists 
	an orthogonal decomposition $ \bigoplus _{d \in D} L_d $ of $ \C^n $ for which:
	\begin{enumerate}
		\item for all $ \lambda \in \R \cap f( \C^n )$, there exists
		a unique $d \in D$ and so a unique space $ L_d $ of the decomposition such that
		 $ f^{-1} ( \{ \lambda \} ) = L_d \setminus \{ \mathbf 0 \}$;
		\item $ f^{-1} ( \{*\} )$ is the residual set 
 \( R = \C^n \setminus \cup _{d \in D} ( L_d \setminus \{ \mathbf 0 \} ) \).
	\end{enumerate}
\end{definition}

\begin{theorem} \label{thm:threeBijections}
	There exist bijections between each pair of the three sets of:
\begin{enumerate}
	\item quantum observables in the form of $n \times n$ Hermitian matrices $ \mathbf A$;
	 \item numerically identified 
	 orthogonal decompositions $ \bigoplus _{ \lambda \in \Lambda } L_\lambda $ of $ \C^n $; 
	\item ortho-measurable functions
	 \( \C^n \owns \mathbf x \mapsto f( \mathbf x) \in \R \cup \{*\} \)
	 satisfying $ f^{-1} ( \{ \lambda \} ) = E_\lambda $ for all $ \lambda \in \Lambda $.
\end{enumerate}
\end{theorem}

\begin{proof}
	First, given any $n \times n$ Hermitian matrix $ \mathbf A$, 
	the spectral theorem \ref{prop:spectralThm} implies the existence 
	of the unique numerically identified orthogonal decomposition 
	 \( \bigoplus\nolimits _{ \lambda \in \Lambda } L_\lambda \) of $ \C^n $,
	 where $ \Lambda $ is the finite spectrum $ s^{ \mathbf A}$ of $ \mathbf A$. 

	Second, let \( \C^n \owns \mathbf x \mapsto f( \mathbf x) \in \R \cup \{*\} \) 
	be any ortho-measurable function.
	Now relabel the sets $ L_d $ 
	in the orthogonal decomposition $ \bigoplus _{d \in D} L_d $ of $ \C^n $ 
	of Definition \ref{def:orthMeasFn} according to the function values $ \lambda $
	in the real part $ \Lambda = \R \cap f( \C^n )$ of the range of $f$, which is a finite set.
	For different values of~$ \lambda $, the corresponding spaces $ L_\lambda $ are orthogonal.
	So the resulting 
	labelled orthogonal decomposition $ \bigoplus _{ \lambda \in \Lambda } L_\lambda $
	is numerically identified.
	Then the decomposition $ \bigoplus _{ \lambda \in \Lambda } L_\lambda $ serves 
	as a parameter of the function $f$ defined by 
\begin{equation} \label{eq:defOMeasFn}
	f \left( \bigoplus\nolimits _{ \lambda \in \Lambda } L_\lambda ; \mathbf x \right) 
	:= \begin{cases} 
	\lambda & \text{if $ \mathbf x \in E_\lambda $}; \\
 	* & \text{if $ \mathbf x \in \C^n \setminus \cup _{ \lambda \in \Lambda } E_\lambda $}.
\end{cases}
\end{equation}
	
	Next, consider the intersection of the graph of this function, 
	which is a subset of the Cartesian product $ \C^n \times ( \R \cup \{*\} )$, 
	with the product set $( \C^n \setminus \{ \mathbf 0\} ) \times \R$. 
	This intersection is the \emph{restricted graph} made up of the finite union 
	of pairwise disjoint sets given by
\begin{equation} \label{eq:graphOMeasFn}
	\Gamma \left( \bigoplus\nolimits _{ \lambda \in \Lambda } L_\lambda \right) 
	:= \bigcup\nolimits _{ \lambda \in \Lambda } \left( E_\lambda \times \{ \lambda \} \right)
\end{equation}
	
	Finally, recalling the spectral decomposition \eqref{eq:spectralDecomp} 
	and putting $ \Lambda = s^{ \mathbf A} $, we have the following chain of bijections
\begin{equation} \label{eq:bijections}
	\mathbf A = \sum\nolimits _{ \lambda \in s^{ \mathbf A} } 
	\lambda \, \boldsymbol \Pi _{ L_\lambda } 
	\longleftrightarrow \bigoplus\nolimits _{ \lambda \in \Lambda } L_\lambda 
	\longleftrightarrow \Gamma \left( \bigoplus\nolimits _{ \lambda \in \Lambda } L_\lambda \right)
\end{equation}
	This chain makes the result evident.
\end{proof}

\end{document}
